%!TEX root = ../deepsmooth.tex

\section{Background and objectives} \label{sec:ivs}

\subsection{The implied volatility surface}

Let $\pi(K, \tau)$ denotes the market price of a call option with time to maturity $\tau>0$ and strike price $K\ge0$.
Without loss of generality, we assume that the initial stock price is $S$, and that the interest-rate $r$ and dividend yield $\delta$ are constants.
Denote by $C(\cdot)$ the Black-Scholes formula, namely
\begin{equation}\label{eq:BSMcp}
C(S, \, \sigma, \, r, \, \delta, \,K, \, \tau) = S^{-\delta \tau} \Phi(d_+) - \e^{-r\tau}K \Phi(d_-),
\end{equation}
where 
$
d_\pm = (\log(S/K) + (r-\delta)\tau)/(\sigma \sqrt{\tau}) \pm (1/2)\sigma\sqrt{\tau}.
$
More details on this formula can be found in Appendix~\ref{sec:app_bs}.
The main objective of this paper is modeling the implied volatility whose definition is given below.
\begin{definition}\label{def:ivs}
The \emph{implied volatility} $\iv(k,\tau)>0$ is given by the equation
\begin{equation} \label{eq:iv_def}
\pi(K,\tau) = C(S, \, \iv(k, \tau), \, r, \, \delta, \, K, \, \tau),
\end{equation}
with the (forward) log moneyness $k=\log(K/S\e^{(r-\delta)\tau})$.
The \emph{implied volatility surface} is given by $\iv(k,\tau)$ for $k\in\R$ and $\tau>0$.
\end{definition}

\textbf{Interpreting the IVS shape.}
For a fixed $\tau>0$,  $\iv(k,\tau)$ with $k\in\R$ defines a volatility smile.
If the smile has a \emph{U shape}, then the tail of the log return $\log(S_T/S_t)$ distribution are thicker than the tails of the Gaussian distribution, and vice versa.
If the smile exhibits a skew, then one side of the log return distribution is thicker than the other.
For example, if the left side of a smile, which is a slice of the surface for a fixed $\tau$, is steeper than the right side, then the log price is more likely to experience large losses than large gains.
The implied volatility surface provides a snapshot representation of valid option prices at a given time point.
Although option prices fluctuate significantly over time, the shape and level of the implied volatility surface is fairly stable and large movements indicate important changes in market conditions.

%In Appendix~\ref{sec:app_bs}, we briefly review the Black-Scholes (BS) option pricing formula which is the baseline model used to extract the implied volatility from market option prices. 
%For more details we refer the reader to standard finance textbooks such as~\cite[Chapter~15]{hull2017options}.

\subsection{Aribtrage-free surface} \label{sec:ao}

A static arbitrage is a static trading strategy that has
a value that is both zero initially and always greater than or equal to zero afterwards,
and a non-zero probability of having a strictly positive value in the future.
In other words, an arbitrage costs nothing to implement while only providing upside potential, that is,
it represents a risk-free investment after accounting for transaction costs.
Under the assumption that economic agents are rational, any such opportunity should be instantaneously exploited until the market is arbitrage free.
Therefore, option pricing models are designed in such a way that their call price surface $\pi(K,T)$ offers no possibility to implement such a strategy.
Standard static arbitrage opportunities are described in Appendix~\ref{sec:staticarb}.

The absence of arbitrage translates into \emph{constraints} on the call price surface $\pi(K,T)$, 
which in turn can be expressed as conditions that the implied volatility surface $\iv(k,\tau)$ must satisfy~\cite{roper2010arbitrage,gatheral2014arbitrage}.
To express those conditions, we define the \emph{total variance} of $\iv(k,\tau)$,
\begin{equation} \label{eq:siv_def}
  \omega(k,\tau)= \iv^2(k,\tau) \, \tau.
\end{equation}

\begin{proposition}{\citet[Theorem~2.9]{roper2010arbitrage}} \label{prop:noAO}
Let $S>0$, $r=\delta=0$, and $\omega:\R \times [0,\infty) \mapsto \R$.
Let $\omega$ satisfy the following conditions:
\begin{enumerate}[label={C\arabic*)}, noitemsep, topsep=0pt]
\item \label{cons:posi} (Positivity) for every $k\in\R$ and $\tau>0$,
$
\omega(k,\tau)>0.
$
\item \label{cons:vmat} (Value at maturity) for every $k\in\R$,
$
\omega(k,0)=0.
$
\item \label{cons:smoo} (Smoothness) for every $\tau>0$, $\omega(\cdot, \, \tau)$ is twice differentiable.
\item \label{cons:mono} (Monotonicity in $\tau$) for every $k\in\R$, $\omega(k, \, \cdot)$ is non-decreasing,
$
\ell_{\rm cal}(k,\tau) = \partial_\tau \omega(k,\tau) \ge 0,
$
where we have written $\partial_\tau$ for $\partial/\partial \tau$.
\item \label{cons:durr} (Durrleman's Condition) for every $\tau>0$ and $k\in\R$,
$$
\ell_{\rm but}(k,\tau) = \left( 1 - \frac{k \, \partial_k \omega(k,\tau)}{2 \omega(k,\tau)}\right)^2 - \frac{\partial_k \omega(k,\tau)}{4} \left(\frac{1}{\omega(k,\tau)} + \frac{1}{4}\right) + \frac{\partial^2_{kk}\omega(k,\tau)}{2} \ge 0,
$$
where we have written $\partial_k$ for $\partial/\partial k$ and $\partial_{kk}$ for $\partial^2/(\partial k\partial k)$
\item \label{cons:asym} (Large moneyness behaviour) for every $\tau>0$, $\sigma^2(k,\tau)$ is linear for $k\rightarrow \pm \infty$.
%$
%\lim_{k\rightarrow\infty} d_+(k, \omega(k,\tau)) = - \infty.
%$
\end{enumerate} 
Then, the resulting call price surface is free of static arbitrage.
\end{proposition}

\labelcref{cons:posi} and~\labelcref{cons:vmat} are necessary conditions that any sensible model must satisfy.
As for~\labelcref{cons:smoo}, it is merely sufficient to prove an absence of arbitrage when~\labelcref{cons:mono},~\labelcref{cons:durr}, and~\labelcref{cons:asym} are also satisfied.
Note that, assuming~\labelcref{cons:smoo}, \labelcref{cons:mono} (respectively~\labelcref{cons:durr} and~\labelcref{cons:asym}) is satisfied if and only if the call price surface is free of calendar spread (butterfly) arbitrage~\cite{gatheral2014arbitrage}.
\labelcref{cons:asym} could be refined by imposing that $\sigma^2(k,\tau)/ \lvert k \rvert < 2$ when $k\rightarrow \pm \infty$ to guarantee the existence of higher order implied moments, see~\cite{lee2004moment,benaim2009regular}.

\begin{remark}
The Proposition~\ref{prop:noAO} is derived under the assumption that $r=\delta=0$ without loss of generality.
Indeed, the static no-arbitrage constraints have the same functional forms as in \ref{cons:mono}--\ref{cons:durr}--\ref{cons:asym} with non-zero parameters and the forward log moneyness $k$ defined in Definition~\ref{def:ivs}. 
\end{remark}

\subsection{Problem formulation}

\textbf{Data availability}.
In practice, market data may need to be validated.
This could be for a variety of reasons.
First, exchange traded securities have at least two quotes, a bid and a ask price, and there is no guarantee that the average prices are arbitrage-free.
Second, the observed prices may not be refreshed and thus not actionable, which may translate into notable input data noise.
In addition, market data is typically sparse away from the money and dense close to the money.
Indeed, far out of the money options are less likely to be exercised and are thus less likely to be used by their buyers.
There is typically more demand and supply for contracts that are around the money.

\textbf{Modeling objectives}.
The goal is to construct an IVS model $\sigma(k,\tau)$ that (I)
generates options prices that are in line with market data,
(II) is free of static arbitrage opportunities in the sense of~\Cref{prop:noAO},
and (III) generalizes to unobserved data regions in a controlled fashion.
